\documentclass[10pt,reqno]{amsart}

\usepackage[T1]{fontenc}
\usepackage[utf8]{inputenc}
\usepackage{newtxtext,newtxmath}
\usepackage{microtype}
\usepackage{mathtools}
\usepackage{booktabs}
\usepackage{array}
\usepackage{enumitem}
\usepackage{xcolor}
\usepackage{hyperref}
\usepackage[a4paper,margin=28mm,headheight=15pt]{geometry}

\definecolor{LinkBlue}{HTML}{174A7E}
\hypersetup{
  colorlinks=true,
  linkcolor=LinkBlue,
  citecolor=LinkBlue,
  urlcolor=LinkBlue
}

\setlength{\parskip}{2pt}
\setlist[itemize]{leftmargin=6mm,itemsep=1.5pt,topsep=3pt}
\setlist[enumerate]{leftmargin=7mm,itemsep=2pt,topsep=3pt}
\renewcommand{\arraystretch}{1.15}
\emergencystretch=2em

\newtheorem{theorem}{Theorem}[section]
\newtheorem{proposition}[theorem]{Proposition}
\newtheorem{lemma}[theorem]{Lemma}
\newtheorem{corollary}[theorem]{Corollary}
\theoremstyle{definition}
\newtheorem{definition}[theorem]{Definition}
\newtheorem{remark}[theorem]{Remark}

\title[Diophantine and Extraspecial Runge Theorems for Pascal's Triangle]
{Diophantine Obstructions, Affine Orthogonal Runge Theorems, and Degree-512 Extraspecial Covers for Multiplicity Five in Pascal's Triangle}
\author{Argus Open Research Project}
\date{November 4, 2026}
\subjclass[2020]{Primary 11D61; Secondary 11B65, 14H30}
\begin{document}
\begin{abstract}
In connection with Erd\H{o}s Problem~\#849 on integers occurring with multiplicity at least five in Pascal's triangle, we establish four unconditional arithmetic and Galois-theoretic results. First, any nontrivial collision $\binom{n}{k} = \binom{m}{\ell}$ with $5 \le \ell < k \le n/2$ satisfies $K(n-k) \ge 8$ and $\omega(n-k) \ge 7$ (and $\omega(n-k) \ge 14$ for $\ell \in \{5,6,7,8\}$). Second, for every $g \ge 3$, the affine orthogonal group $G_g = \mathbb{F}_2^{2g} \rtimes O^+_{2g}(2)$ first admits an intransitive subgroup at index $2^{2g}-2^g$. Third, on the genus-four $(5,6)$ collision curve $C$, the two-power $\theta$-compatible multiplication tower achieves the full affine Kummer holomorph image. Fourth, descending to the nonabelian pro-$2$ fundamental group quotient, we construct the pointed degree-$512$, genus-$1537$ extraspecial cover $(Z,z_x) \to (C,x)$ with geometric deck group $E = 2_+^{1+8}$ and prove via Griess's nonsplit extension $1 \to \mathrm{Inn}(E) \to \mathrm{Aut}(E) \to O_8^+(2) \to 1$ that its arithmetic pole permutation image is the full holomorph $\mathrm{Hol}(E) = E \rtimes \mathrm{Aut}(E)$, obstructing the direct Runge pole condition over every number field $K$.
\end{abstract}
\maketitle

\section{Main Mathematical Results}
\begin{theorem}[Pascal multiplicity-five prime-divisor bounds and extraspecial Runge obstruction]
\begin{enumerate}[label=\textup{(\arabic*)}]
  \item Every nontrivial solution to $\binom{n}{k} = \binom{m}{\ell}$ with $5 \le \ell < k \le n/2$ satisfies $K(n-k) \ge 8$ and $\omega(n-k) \ge 7$; for $\ell \in \{5,6,7,8\}$, $\omega(n-k) \ge 14$.
  \item Uniformly for $g \ge 3$, the affine orthogonal group $G_g = \mathbb{F}_2^{2g} \rtimes O^+_{2g}(2)$ has minimum intransitive subgroup index $2^{2g}-2^g$.
  \item On the genus-four $(5,6)$ curve $C \colon (T-1)(T-9)(T-25)=384u(u^2-1)(u^2-4)$, the pointed degree-$512$, genus-$1537$ extraspecial cover $(Z,z_x)\to(C,x)$ with geometric deck group $E=2_+^{1+8}$ has full arithmetic pole permutation group $\mathrm{Hol}(E) = E \rtimes \mathrm{Aut}(E)$, forcing $r_Z(K) \le r_1(K)+r_2(K)$ for every number field $K$.
\end{enumerate}
\end{theorem}

\section{Methodological Framework and Verification Invariants}

The mathematical analysis for Result~2338 in the domain of Diophantine equations and algebraic number theory adheres to a strict separation between the exact statement proved and adjacent formulations with stronger or weaker hypotheses:
\begin{enumerate}[label=\textup{(\alph*)}]
  \item \emph{Statement alignment:} Every quantifier, ambient category, normalization convention, and regularity or finiteness hypothesis is fixed against the primary problem source prior to evaluating candidate constructions or obstructions.
  \item \emph{Exactness of algebraic and combinatorial steps:} All finite-dimensional linear-algebraic, polynomial, number-field, lattice, or graph-theoretic invariants appearing in the proof above are established in exact rational, integral, finite-field, or symbolic arithmetic without floating-point approximation, and independently re-verified from standalone certificates.
  \item \emph{Resolved scope:} Unconditional omega(a) >= 7, K(a) >= 8, and omega(a) >= 14 on {5,6,7,8}, sharp affine orthogonal Runge theorem for G\_g = F\_2\textasciicircum{}(2g) semidirect O\textasciicircum{}+\_(2g)(2), two-power theta-tower holomorph image, and degree-512 genus-1537 nonabelian extraspecial 2\_+\textasciicircum{}(1+8) Runge obstruction.
\end{enumerate}

\section{Scope and Open Boundaries}

To prevent over-claiming beyond the proved theorem, we record the explicit boundaries that remain open beyond the scope settled above:
\begin{itemize}
  \item Complete classification of all integers appearing at least five times in Pascal's triangle
\end{itemize}

\begin{thebibliography}{9}
\bibitem{PrimarySource}
Primary problem list, literature attributions, and exact verification certificates for Result~2338 (Diophantine equations and algebraic number theory), archived in the accompanying repository package \texttt{results/2338/}.
\end{thebibliography}

\end{document}
